% Завдання до курсової роботи — за Додатком Б навчального посібника.
\structural{ЗАВДАННЯ}

Виконати дослідження для такої системи.

На обробку до інформаційної системи приймаються три типи завдань A, B і C. Виходячи з наявності
оперативної памʼяті обчислювальної системи, завдання типів A і B можуть виконуватися одночасно. Тобто
завдання типу A можуть виконуватися паралельно із завданням свого класу або із завданням типу B.
Аналогічно, завдання типу B можуть виконуватися паралельно із завданням свого класу або із завданням
типу A. Завдання типу C монополізують інформаційну систему. Завдання надходять у середньому з
інтенсивністю 0,2~завдання/хв, 0,07~завдання/хв, 0,06~завдання/хв для класів A, B та C відповідно.
Завдання виконуються протягом часу, що є випадковою величиною, розподіленою за експоненціальним законом
з інтенсивністю 0,25~завдання/хв, 0,2~завдання/хв, 0,1~завдання/хв для класів A, B та C відповідно.

Упорядкування в черзі на обслуговування визначається комбінацією пріоритетів завдань. Можливі такі
комбінації:
\begin{itemize}
\item A -- вищий, B -- середній, C -- низький;
\item B -- вищий, A -- середній, C -- низький;
\item C -- вищий, B -- середній, A -- низький;
\item C -- вищий, A -- середній, B -- низький;
\item C -- вищий, A і B -- низький.
\end{itemize}

Метою моделювання є визначення впливу різних стратегій упорядкування в черзі при різних значеннях
інтенсивності надходження на такі параметри інформаційної системи: середній час обслуговування завдання
в системі; середня кількість завдань, що очікують обслуговування; середній час очікування в чергах.

Для вищезазначеної системи:
\begin{enumerate}
\item розробити опис концептуальної моделі системи;
\item виконати формалізацію опису обʼєкта моделювання в термінах стохастичної мережі Петрі;
\item розробити алгоритм імітації моделі дискретно-подійної системи у відповідності до побудованого
  формального опису;
\item для доведення коректності побудованого алгоритму виконати верифікацію алгоритму імітації;
\item визначити статистичні оцінки заданих характеристик моделі, що є метою моделювання;
\item провести експериментальне дослідження моделі;
\item інтерпретувати результати моделювання та сформулювати пропозиції щодо поліпшення функціонування
  системи;
\item зробити висновки щодо складності розробки моделі та алгоритму імітації на основі використаного
  формалізму, отриманих результатів моделювання та їх корисності.
\end{enumerate}
